\section{Challenge Tasks related to $p(z)$ estimation}\label{sec:tasks}
\label{sec:pz_tasks}

\subsection{Task set 1: Estimate redshifts using representative
  training samples}

The first, simplest task is to estimate redshifts using
representative training samples. I.e., the training samples are drawn
from the same distributions as the test samples. For this task set we
did not use any of the spectroscopic selection emulation, but simply
applied a uniform magnitude cut of $i < 23$ in selecting objects for both
the training and test samples.

The four 
\texttt{pz\_challenge\_taskset\_1\_\{simulation\}\_training\_\{scenario\}.hdf5}
files are the training sets for the ``Flagship'' and ``Cardinal''
simulations, emulating 1 year and 10 years of LSST data under the
expected observing strategy and conditions. These files have true
redshifts to serve as labels.

The corresponding 
\texttt{pz\_challenge\_taskset\_1\_\{simulation\}\_test\_\{scenario\}.hdf5}
files were drawn from the same distributions. The true redshifts have
been removed from these files. The task is to assign $p(z)$
estimates for all the objects in these 4 test files.

The subtasks in this task set are:

\begin{enumerate}
  \item{Estimate $p(z)$ for each object in each of the test files and
      provide the estimates in a downloadable \texttt{tar} file.}
  \item{Provide pre-trained models appropriate to each of the training
      files and implement a Python function (\texttt{run\_taskset\_1\_estimation\_only}) to use those
      pre-trained models to estimate $p(z)$ for each object in the
      associated test files.}    
  \item{Implement a Python function (\texttt{run\_taskset\_1\_training\_and\_estimation}) to train a model for
      each training file and use that model to estimate $p(z)$ for
      each object in the associated test files.}    
\end{enumerate}


\subsection{Task set 2: Estimate redshifts on non-representative
  samples}

The second, slightly more challenging task is to estimate redshifts
using non-representative training samples. I.e., the training samples
are not drawn from the same distributions as the test samples. For this task set we
applied the spectroscopic selection emulation for the training set, but
retained all the objects down to $i < 25.4$ in the test set.
Accordingly, the training set will not be representative of the
fainter objects in the test set. This reflects that spectroscopic redshifts
are typically significantly more difficult to obtain than photometry.

The four 
\texttt{pz\_challenge\_taskset\_2\_\{simulation\}\_training\_\{scenario\}.hdf5}
files are the training sets for the ``Flagship'' and ``Cardinal''
simulations, emulating 1 year and 10 years of LSST data under the
expected observing strategy and conditions and with spectroscopic
selections emulated.

The corresponding 
\texttt{pz\_challenge\_taskset\_2\_\{simulation\}\_test\_\{scenario\}.hdf5}
files were drawn from the distributions of all objects down to $i <
25.4$, and the true redshifts have been removed from these files. 
The task is to assign $p(z)$ estimates for all the objects in these 4 test files.

The subtasks in this task set are:

\begin{enumerate}
  \item{Estimate $p(z)$ for each object in each of the test files and
      provide the estimates in a downloadable \texttt{tar} file.}
  \item{Provide pre-trained models appropriate to each of the training
      files and implement a Python function
      (\texttt{run\_taskset\_2\_estimation\_only}) to use those
      pre-trained models to estimate $p(z)$ for each object in the
      associated test files.}    
  \item{Implement a Python function
      (\texttt{run\_taskset\_2\_training\_and\_estimation}) to train a model for
      each training file and use that model to estimate $p(z)$ for
      each object in the associated test files.}    
\end{enumerate}
  
